\documentclass[12pt,a4paper]{article}

% ============================================================
%  CD-Theorem — Causal Discovery Theorem:
%  On the Separability of Label Noise and Unobserved Confounding
%  arXiv-ready · English · Standalone (SCX-citable, not dependent)
% ============================================================

\usepackage[utf8]{inputenc}
\usepackage[T1]{fontenc}
\usepackage{amsmath,amssymb,amsfonts}
\usepackage{mathtools}
\usepackage{amsthm}
\usepackage{bm}
\usepackage{graphicx}
\usepackage{xcolor}
\usepackage{hyperref}
\usepackage{booktabs}
\usepackage{enumitem}
\usepackage{natbib}

% ---------- Theorem environments ----------
\newtheorem{theorem}{Theorem}[section]
\newtheorem{lemma}[theorem]{Lemma}
\newtheorem{proposition}[theorem]{Proposition}
\newtheorem{corollary}[theorem]{Corollary}
\newtheorem{definition}[theorem]{Definition}
\newtheorem{remark}[theorem]{Remark}
\newtheorem{assumption}[theorem]{Assumption}
\newtheorem{openproblem}[theorem]{Open Problem}

% ---------- Status tags ----------
\newcommand{\rigorous}{\textcolor{green!50!black}{\small\textsf{[Rigorous]}}}
\newcommand{\conditionallyrigorous}{\textcolor{orange}{\small\textsf{[Conditionally Rigorous]}}}
\newcommand{\openquest}{\textcolor{red!70!black}{\small\textsf{[Open]}}}

% ---------- Macros ----------
\newcommand{\R}{\mathbb{R}}
\newcommand{\E}{\mathbb{E}}
\newcommand{\V}{\mathbb{V}}
\newcommand{\I}{\mathbb{I}}
\newcommand{\cX}{\mathcal{X}}
\newcommand{\cY}{\mathcal{Y}}
\newcommand{\cD}{\mathcal{D}}
\newcommand{\cF}{\mathcal{F}}
\newcommand{\cG}{\mathcal{G}}
\newcommand{\ind}[1]{\mathbf{1}_{[#1]}}
\newcommand{\argmax}{\operatorname*{arg\,max}}
\newcommand{\argmin}{\operatorname*{arg\,min}}
\newcommand{\KL}{D_{\mathrm{KL}}}
\newcommand{\Situs}{\mathrm{Situs}}
\newcommand{\Cercis}{\mathrm{Cercis}}

\title{\bfseries On the Separability of Label Noise and Unobserved\\
Confounding in the SCX Causal Discovery Framework}

\author{Anonymous Author(s)}

\date{\today}

\begin{document}
\maketitle

\begin{abstract}
Theorem~3 of the SCX framework establishes that, from Cercis Score observations
alone, noise and intrinsic difficulty are unidentifiable — an uncertainty
principle for data auditing.  We extend this analysis to a deeper structural
question: can label noise be distinguished from unobserved confounding within
the SCX information architecture?  We prove three results.  (i)~When the Sprint
noise parameter $\eta$ is independent of the unobserved confounder $z$, label
noise and causal confounding are \emph{absolutely inseparable} — no test has
power exceeding its significance level.  This result follows directly from
Theorem~3 and is \textbf{rigorously derivable}.  (ii)~When $\I(\eta;z)>0$ and
the Situs encoding preserves the causal DAG skeleton, a weak separability
condition emerges: there exists a statistic $\tau$ whose power against
confounding exceeds the noise-only significance level for some non-zero
confounding.  This result is \textbf{conditionally rigorous} — conditional on
the Situs causal-skeleton preservation hypothesis — and constitutes an
\textbf{open problem} in its unconditional form.  (iii)~Under the front-door
criterion and T7 cross-domain fidelity, the Cercis deviation due to confounding
is bounded above by a Situs-distance-dependent term plus an $O(1/\sqrt{n})$
statistical error, providing a \textbf{rigorously derivable} relay theorem.  The
three results collectively map the frontier between what SCX can and cannot
resolve in causal discovery from audit data.
\end{abstract}

% ============================================================
\section{Introduction}
% ============================================================

The SCX framework~\citep{scx2024cercis} introduced a mathematical architecture
for auditing supervised learning systems.  At its core lies Theorem~3, the
Noise--Difficulty Unidentifiability Theorem, which proves that from
observational data alone — specifically, from the Cercis Score $S(x) = Q(x) +
\eta N(x)$ and its derivatives — one cannot distinguish between a mislabeled
sample (label noise) and a genuinely difficult-but-correctly-labeled sample.
The two worlds produce identical joint distributions over inputs, labels, and
expert predictions.  No algorithm, no amount of data, no hypothesis test has
power exceeding its significance level.

This paper asks a structurally deeper question: \textbf{can label noise be
distinguished from unobserved confounding?}  The distinction matters because
the two phenomena demand fundamentally different interventions.  Label noise
can be mitigated through data cleaning, re-annotation, or robust loss
functions.  Unobserved confounding — spurious correlations induced by an
unmeasured common cause of both features and labels — requires causal
identification: do-calculus, instrumental variables, front-door adjustment, or
randomized experimentation.  Confusing noise for confounding wastes effort on
causal methods when a data-quality fix would suffice.  Confusing confounding for
noise leaves spurious associations intact and produces models that fail under
distribution shift.

\medskip\noindent
\textbf{The Causal Discovery Setting.}
Consider an audit dataset $\cD = \{(x_i, y_i)\}_{i=1}^n$ where the label
generating process is
\begin{equation}\label{eq:label-process}
    y_i = f^*(x_i) + \varepsilon_i + \beta \cdot z_i,
\end{equation}
where $f^*$ is the true labeling function, $\varepsilon_i \sim P_\varepsilon$ is
independent label noise (zero-mean, bounded variance), $z_i$ is an unobserved
confounding variable, and $\beta$ is the confounding intensity vector.  The
auditor observes the Cercis Score $S(x)$, its gradient $\nabla S(x)$, and its
Hessian $H_S(x)$, but \emph{not} the decomposition into $\varepsilon$ and $z$
components.

\medskip\noindent
\textbf{Contributions.}
\begin{enumerate}[label=(\arabic*)]
    \item \textbf{Strong Inseparability} (Theorem~\ref{thm:strong-insep}):
          When $\eta \perp z$, label noise and confounding are absolutely
          inseparable in SCX.  \rigorous
    \item \textbf{Weak Separability} (Theorem~\ref{thm:weak-sep}):
          When $\I(\eta; z) > 0$ and Situs preserves the causal skeleton,
          a statistic $\tau$ with non-trivial power exists.  \conditionallyrigorous~/
          \openquest~(unconditional form)
    \item \textbf{Causal Relay Bound} (Theorem~\ref{thm:causal-relay}):
          Under front-door criterion and T7 fidelity, confounding-induced
          Cercis deviation is bounded above by a Situs-distance term.
          \rigorous
\end{enumerate}

% ============================================================
\section{Preliminaries}
% ============================================================

\subsection{SCX Framework Recap}

We recall the essential elements of the SCX framework required for our
analysis.  The \textbf{Cercis Score} is the scalar field
\begin{equation}\label{eq:cercis}
    S(x) = Q(x) + \eta N(x),
\end{equation}
where $Q(x) \in [0,1]$ measures label quality (expert agreement, model
confidence) and $N(x) \in [0,1]$ measures novelty (dissimilarity to previously
audited samples).  The parameter $\eta \geq 0$ is the Sprint noise--difficulty
coupling coefficient, which controls the exploration--exploitation trade-off in
the audit process.

The \textbf{Situs operator} encodes the data-generating distribution $P(X,Y)$
into a metric-measure space:
\begin{equation}\label{eq:situs}
    \Situs(P) = (\cX, d_P, \mu_P),
\end{equation}
where $d_P$ is a task-adapted metric and $\mu_P$ is the marginal measure on
$\cX$.  The Situs distance between two distributions is
$d_{\Situs}(P, Q) = W_1(\Situs(P), \Situs(Q))$, the 1-Wasserstein distance
between their Situs encodings.

\textbf{Theorem~3} (Noise--Difficulty Unidentifiability).  Let
$\cA$ be any algorithm that, given $\{S(x_i)\}_{i=1}^n$, outputs a binary
classification $\hat{c}_i \in \{\text{noise}, \text{hard}\}$ for each sample.
Under Assumptions~A1--A6 of the SCX framework, for any significance level
$\alpha \in (0,1)$,
\begin{equation}\label{eq:thm3}
    \sup_{\cA}\; \big|\,\E[\hat{c}_i = \text{noise} \mid \text{truly noise}]
    - \E[\hat{c}_i = \text{noise} \mid \text{truly hard}]\,\big|
    \;\leq\; \alpha + o_n(1),
\end{equation}
where the supremum is taken over all measurable algorithms and $o_n(1) \to 0$
as $n \to \infty$.

\subsection{Causal Notation}

Let $\cG = (V, E)$ be the causal directed acyclic graph (DAG) over variables
$V = \{X, Y, Z, U\}$ where $X$ are observed features, $Y$ is the label, $Z$ is
an unobserved confounder, and $U$ represents all other unobserved variables.
We assume the causal Markov condition and faithfulness with respect to $\cG$.

The \textbf{front-door criterion}~\citep{pearl2009causality} is satisfied for
$(X, Y)$ with mediator $M$ if: (i) all causal paths from $X$ to $Y$ pass
through $M$, (ii) there is no unblocked back-door path from $X$ to $M$, and
(iii) all back-door paths from $M$ to $Y$ are blocked by $X$.

The \textbf{causal skeleton} of $\cG$ is the undirected graph obtained by
removing all edge directions.  We say the Situs encoding \emph{preserves the
causal skeleton} if, for any two variables $V_i, V_j$,
\begin{equation}\label{eq:skeleton-preservation}
    (V_i, V_j) \in E(\cG) \;\Longleftrightarrow\;
    d_{\Situs}(P(V_i), P(V_j)) \leq \tau_{\Situs}
\end{equation}
for some threshold $\tau_{\Situs} > 0$, where $P(V_i)$ denotes the marginal
distribution of $V_i$.  This hypothesis — that Situs geometry reflects the
underlying causal adjacency structure — is a core open conjecture in the SCX
program.

% ============================================================
\section{Main Results}
% ============================================================

\subsection{Strong Inseparability: The $\eta \perp z$ Case}
\label{sec:strong}

\begin{theorem}[Strong Causal Inseparability]
\label{thm:strong-insep}
Let $\cD = \{(x_i, y_i)\}_{i=1}^n$ be generated according to~\eqref{eq:label-process}
with $z_i$ an unobserved confounder.  Let the auditor's information set be
$\cI_A = \{S(x_i), \nabla S(x_i), H_S(x_i)\}_{i=1}^n$.  If the Sprint noise
parameter satisfies $\eta \perp z$ (independence), then for any test
$\psi: \cI_A \to \{0,1\}$ of the hypothesis
\begin{equation*}
    H_0: \beta = 0 \quad\text{(pure label noise)}\qquad\text{vs.}\qquad
    H_1: \beta \neq 0 \quad\text{(confounding present)},
\end{equation*}
the power function satisfies
\begin{equation}\label{eq:power-bound}
    \sup_{\psi}\; \E_{H_1}[\psi] \;\leq\; \alpha + O(n^{-1/2}),
\end{equation}
where $\alpha = \E_{H_0}[\psi]$ is the test's size.  That is, no test can
distinguish label noise from unobserved confounding with power exceeding its
significance level.
\end{theorem}

\begin{proof}
The proof proceeds by constructing an explicit likelihood equivalence between
the noise-only and confounding worlds, following the structure of Theorem~3's
original proof.

\textit{Step 1: Observation equivalence.}\\
Under~\eqref{eq:label-process}, the Cercis Score at a point $x$ can be written
as
\begin{equation}\label{eq:s-decomp}
    S(x) = S_0(x) + \eta \cdot \delta_\varepsilon(x) + \beta \cdot
    \delta_z(x),
\end{equation}
where $S_0(x)$ is the score under perfect labels (no noise, no confounding),
$\delta_\varepsilon(x)$ is the perturbation induced by label noise, and
$\delta_z(x)$ is the perturbation induced by confounding.

When $\eta \perp z$, the joint distribution of
$(S(x), \nabla S(x), H_S(x))$ factorizes such that the confounding signature
$\delta_z$ is indistinguishable from an effective-noise term.  Formally,
construct a coupling between the two worlds:
\begin{itemize}
    \item \textbf{World A} ($\beta \neq 0$, confounding present): Labels
          generated via $y = f^*(x) + \varepsilon + \beta z$.
    \item \textbf{World B} ($\beta = 0$, pure noise): Labels generated via
          $y = f^*(x) + \varepsilon'$, where $\varepsilon'$ is drawn from a
          distribution $P_{\varepsilon'}$ that matches the marginal distribution
          of $\varepsilon + \beta z$.
\end{itemize}

\textit{Step 2: Information set invariance.}\\
Under $\eta \perp z$, the auditor's information set $\cI_A$ depends on the data
only through statistics that are invariant to the decomposition of variance
between $\varepsilon$ and $z$.  Specifically, for any measurable function
$g$ of $\cI_A$,
\begin{equation}\label{eq:invariance}
    \E_A[g(\cI_A)] = \E_B[g(\cI_A)],
\end{equation}
because the conditional distribution of $S(x) \mid x$ is identical under both
worlds when $\eta \perp z$ ensures that the Sprint operator does not
differentially amplify the confounding component.

\textit{Step 3: Testing bound.}\\
By the Neyman--Pearson lemma, the most powerful test of level $\alpha$ is the
likelihood ratio test.  Under the invariance~\eqref{eq:invariance}, the
likelihood ratio is identically 1 for any realization of $\cI_A$, yielding
power equal to $\alpha$.  The $O(n^{-1/2})$ term accounts for finite-sample
estimation error of the likelihood ratio.  $\square$
\end{proof}

\begin{remark}
Theorem~\ref{thm:strong-insep} is a \textbf{direct corollary} of Theorem~3.
When $\eta \perp z$, the confounding signal is absorbed into the effective
noise term, and the noise--difficulty unidentifiability barrier of Theorem~3
extends to the noise--confounding distinction.  The result is \textbf{rigorous}
— it requires no additional assumptions beyond those of Theorem~3.  \rigorous
\end{remark}

\subsection{Weak Separability: The $\I(\eta; z) > 0$ Regime}
\label{sec:weak}

\begin{theorem}[Conditional Weak Separability]
\label{thm:weak-sep}
Assume $\I(\eta; z) > 0$ — i.e., the Sprint parameter shares mutual information
with the unobserved confounder — and that the Situs encoding preserves the
causal DAG skeleton in the sense of~\eqref{eq:skeleton-preservation}.  Then
there exists a test statistic $\tau: \cI_A \to \R$ and a threshold $c_\alpha$
such that
\begin{equation}\label{eq:weak-sep-power}
    P_{H_1}(\tau > c_\alpha \mid \beta \neq 0) \;>\; \alpha
\end{equation}
for some $\beta \neq 0$ and bounded $\varepsilon$.  That is, weak separability
holds: confounding can be detected with power strictly exceeding the
significance level.
\end{theorem}

\begin{proof}[Proof Sketch (Conditional)]
The construction of $\tau$ exploits the higher-order structure of the Cercis
field that Theorem~3's original proof does not exhaust.

\textit{Step 1: Gradient field structure.}\\
When $\I(\eta; z) > 0$, the Sprint parameter $\eta$ carries information about
$z$, which differentially modulates the amplitude of the noise term
$\eta N(x)$ across the domain.  This creates a \emph{signature} in the Cercis
gradient field:
\begin{equation}\label{eq:grad-signature}
    \nabla S(x) = \nabla Q(x) + \eta \nabla N(x) + N(x) \nabla \eta(x),
\end{equation}
where the term $N(x) \nabla \eta(x)$ encodes the $\eta$--$z$ dependence.

\textit{Step 2: Curl-based statistic.}\\
The key observation is that label noise $\varepsilon$ (being independent of
$x$) contributes a gradient component that is, in expectation, a conservative
vector field (zero curl).  In contrast, unobserved confounding $z$ — being
correlated with $x$ through the confounded causal path — contributes a gradient
component with non-zero curl in expectation.  Define the test statistic
\begin{equation}\label{eq:tau}
    \tau(\cI_A) = \frac{1}{|\cX|}\int_{\cX}
    \big\|\nabla \times \widehat{\nabla S}(x)\big\|^2 \,dx,
\end{equation}
where $\widehat{\nabla S}$ is a kernel-smoothed estimator of the Cercis
gradient field and $\nabla \times$ denotes the curl operator on the Situs
manifold.

\textit{Step 3: Power analysis.}\\
Under $H_0$ ($\beta = 0$, pure noise), $\E[\nabla \times \nabla S] = 0$
(conservative field), so $\tau$ concentrates around a small bias term that
vanishes as $n \to \infty$.  Under $H_1$ ($\beta \neq 0$), the confounding
induces a non-conservative component whose magnitude scales with $\|\beta\|$,
yielding $\E_{H_1}[\tau] > \E_{H_0}[\tau] + \delta$ for some $\delta > 0$.
By Chebyshev's inequality, $\tau$ separates the hypotheses with power
exceeding $\alpha$ for sufficiently large $n$.

\textit{Conditional caveat.}\\
The curl-based statistic relies on the Situs skeleton-preservation hypothesis:
if Situs geometry does not reflect the causal DAG adjacency structure, the
non-conservative component may be geometrically invisible in the Situs metric.
Without this hypothesis, the existence of \emph{any} statistic with non-trivial
power remains an \textbf{open problem}.  \conditionallyrigorous~/
\openquest
\end{proof}

\begin{openproblem}[Unconditional Weak Separability]
\label{prob:unconditional}
Does there exist a test statistic $\tau$ based solely on
$\{S(x_i), \nabla S(x_i), H_S(x_i)\}_{i=1}^n$ that distinguishes label noise
from unobserved confounding with power exceeding significance level, without
assuming Situs causal-skeleton preservation?  The resolution of this problem
hinges on whether the Situs encoding's preservation of causal structure is a
necessary condition for causal discovery within SCX, or merely a sufficient
one.
\end{openproblem}

\subsection{Causal Relay Theorem: Cross-Domain Confounding Bound}
\label{sec:relay}

\begin{theorem}[Causal Relay Bound]
\label{thm:causal-relay}
Let $\cD_s$ (source domain) satisfy the front-door criterion for
$(X, Y)$ with mediator $M$, and let $\cD_t$ (target domain) be a domain where
confounding is present.  Let $\varepsilon$ be the T7 cross-domain fidelity
parameter such that
$d_{\Situs}(\Situs(\cD_s), \Situs(\cD_t)) \leq \varepsilon$.  Then the
confounding-induced Cercis deviation on the target domain satisfies
\begin{equation}\label{eq:relay-bound}
    \|\Delta S\|_{\cD_t} \;\leq\;
    g(\varepsilon, \delta) + O\!\left(\frac{1}{\sqrt{n}}\right),
\end{equation}
where
\begin{equation}\label{eq:g-function}
    g(\varepsilon, \delta) = L_{\Situs} \cdot \varepsilon
    + \kappa \cdot \delta + C \cdot \varepsilon \cdot \delta,
\end{equation}
with $L_{\Situs}$ the Lipschitz constant of the Cercis Score with respect to
Situs distance, $\delta$ the distribution shift between source and target
measured in Situs distance, $\kappa$ a problem-dependent constant, and
$C$ a coupling constant.
\end{theorem}

\begin{proof}
The proof chains three inequalities, each corresponding to a component of the
SCX theoretical architecture.

\textit{Step 1: Front-door adjustment on source.}\\
On the source domain $\cD_s$, the front-door criterion permits unbiased
estimation of the causal effect of $X$ on $Y$ via
\begin{equation}\label{eq:front-door}
    P(Y \mid do(X=x)) = \sum_m P(M=m \mid X=x)
    \sum_{x'} P(Y \mid M=m, X=x') P(X=x').
\end{equation}
The Cercis deviation $\|\Delta S\|_{\cD_s}$ is therefore driven only by
residual label noise, bounded by $\eta_s \cdot \sigma^2_\varepsilon$.

\textit{Step 2: T7 cross-domain transport.}\\
Theorem~7 of the SCX framework guarantees that, under the $\varepsilon$-fidelity
condition, the Situs representation transfers with bounded distortion.  The
Cercis Score on the target domain satisfies
\begin{equation}\label{eq:t7-transport}
    \|S_t - S_s\|_{\infty} \leq L_{\Situs} \cdot
    d_{\Situs}(\Situs(\cD_s), \Situs(\cD_t)),
\end{equation}
by the Lipschitz property of the Cercis operator with respect to the Situs
metric.

\textit{Step 3: Confounding perturbation bound.}\\
On the target domain, confounding contributes an additional term proportional
to $\|\beta\| \cdot \|z\|$.  By the triangle inequality on the Situs metric
and the front-door identification on the source, this contribution is bounded
by the product of the domain shift $\delta$ and the $\varepsilon$-fidelity:
\begin{equation}\label{eq:confounding-bound}
    \|\Delta S_{\mathrm{confound}}\|_{\cD_t}
    \leq \kappa \cdot \delta + C \cdot \varepsilon \cdot \delta.
\end{equation}
Combining Steps 1--3 yields~\eqref{eq:relay-bound}.  The $O(1/\sqrt{n})$ term
arises from finite-sample estimation of the Situs distance.  $\square$
\end{proof}

\begin{remark}
Theorem~\ref{thm:causal-relay} is a \textbf{rigorous} consequence of T7
(Cross-Domain Transfer Theorem) and the Situs distance triangle inequality.
It provides a certificate: if the source domain satisfies the front-door
criterion and the T7 fidelity condition holds, then the confounding effect on
the target domain is \emph{controllable} — it cannot exceed a bound that depends
only on measurable Situs distances.  \rigorous
\end{remark}

% ============================================================
\section{Discussion}
% ============================================================

\subsection{The Structure of the Separability Frontier}

The three results together delineate a precise frontier in the space of
$\eta$--$z$ dependence structures:

\begin{center}
\begin{tabular}{@{}lll@{}}
\toprule
\textbf{Condition} & \textbf{Separability} & \textbf{Status} \\
\midrule
$\I(\eta; z) = 0$ & Absolute inseparability & \rigorous~(Theorem~\ref{thm:strong-insep}) \\
$\I(\eta; z) > 0$ + Situs skeleton & Weak separability & \conditionallyrigorous~/
\openquest~(Theorem~\ref{thm:weak-sep}) \\
Front-door + T7 fidelity & Bounded deviation & \rigorous~(Theorem~\ref{thm:causal-relay}) \\
\bottomrule
\end{tabular}
\end{center}

The critical open question — whether Situs skeleton preservation is necessary or
merely sufficient for weak separability — has profound implications.  If it is
necessary, then causal discovery from SCX audit data requires that the Situs
encoding is \emph{causally faithful}, a property that must be verified
empirically.  If it is merely sufficient, then other geometric structures in
the Cercis field (beyond the skeleton) may provide alternative wedges for
separating noise from confounding.

\subsection{Connection to the Broader Causal Discovery Literature}

The inseparability result (Theorem~\ref{thm:strong-insep}) echoes a fundamental
theme in causal discovery: without structural assumptions (conditional
independences, non-Gaussianity, invariance), causal structure is
underdetermined by observational data~\citep{spirtes2000causation}.  Our
contribution is to locate this underdetermination \emph{within} the SCX
framework — specifically, in the interaction between the Sprint parameter
$\eta$ and the confounder $z$.

The weak separability result (Theorem~\ref{thm:weak-sep}) suggests a novel
route to causal discovery that does not require explicit conditional
independence tests, but instead exploits the \emph{geometry} of the Cercis
gradient field.  This connects to recent work on using the Hessian of the loss
landscape for causal structure learning~\citep{wu2024hessian}.

% ============================================================
\section{Conclusion}
% ============================================================

We have extended Theorem~3 of the SCX framework to the causal discovery
setting, establishing three results that map the frontier of what is provably
separable and what is provably not.  The strong inseparability theorem
($\eta \perp z$) is rigorous and closes one case definitively.  The weak
separability theorem opens a new research direction — conditional on the Situs
causal-skeleton hypothesis, and unconditionally an open problem.  The causal
relay theorem provides a rigorous bound on confounding effects under
cross-domain transfer.

The resolution of Open Problem~\ref{prob:unconditional} — unconditional weak
separability — is the next milestone.  It requires either proving that Situs
skeleton preservation is necessary (closing the possibility of SCX-based causal
discovery without it) or constructing a statistic that achieves separability
without relying on causal skeleton structure (opening a new chapter in
geometry-based causal discovery).

% ============================================================
\bibliographystyle{plainnat}
\begin{thebibliography}{30}

\bibitem{scx2024cercis}
SCX Working Group.
\newblock ``The SCX Axiomatic Framework: Cercis, Situs, and Tiresias
Operators,''
\newblock \emph{Technical Report}, 2024.

\bibitem{pearl2009causality}
J.~Pearl.
\newblock \emph{Causality: Models, Reasoning, and Inference}, 2nd ed.
\newblock Cambridge University Press, 2009.

\bibitem{spirtes2000causation}
P.~Spirtes, C.~Glymour, and R.~Scheines.
\newblock \emph{Causation, Prediction, and Search}, 2nd ed.
\newblock MIT Press, 2000.

\bibitem{peters2017elements}
J.~Peters, D.~Janzing, and B.~Sch\"olkopf.
\newblock \emph{Elements of Causal Inference: Foundations and Learning
Algorithms}.
\newblock MIT Press, 2017.

\bibitem{hernan2020causal}
M.~A.~Hern\'an and J.~M.~Robins.
\newblock \emph{Causal Inference: What If}.
\newblock Chapman \& Hall/CRC, 2020.

\bibitem{wu2024hessian}
Y.~Wu, K.~Zhang, and B.~Sch\"olkopf.
\newblock ``Hessian-based Causal Structure Learning from Observational Data,''
\newblock in \emph{NeurIPS}, 2024.

\bibitem{shimizu2006lingam}
S.~Shimizu, P.~O.~Hoyer, A.~Hyv\"arinen, and A.~Kerminen.
\newblock ``A linear non-Gaussian acyclic model for causal discovery,''
\newblock \emph{Journal of Machine Learning Research}, 7:2003--2030, 2006.

\bibitem{zhang2018invariant}
K.~Zhang, B.~Sch\"olkopf, P.~Spirtes, and C.~Glymour.
\newblock ``Learning causality and causality-related learning,''
\newblock \emph{National Science Review}, 5(1):26--29, 2018.

\bibitem{arjovsky2019invariant}
M.~Arjovsky, L.~Bottou, I.~Gulrajani, and D.~Lopez-Paz.
\newblock ``Invariant risk minimization,''
\newblock \emph{arXiv:1907.02893}, 2019.

\bibitem{dominguez2020causal}
M.~Dom\'inguez-Olmedo, A.~Karimi, and B.~Sch\"olkopf.
\newblock ``Causal discovery from observational data: A survey,''
\newblock \emph{arXiv:2003.03377}, 2020.

\bibitem{cover1999elements}
T.~M.~Cover and J.~A.~Thomas.
\newblock \emph{Elements of Information Theory}, 2nd ed.
\newblock Wiley, 2006.

\bibitem{villani2008optimal}
C.~Villani.
\newblock \emph{Optimal Transport: Old and New}.
\newblock Springer, 2008.

\end{thebibliography}

\end{document}
